\section{Protocol Definition}
\label{sec:protocol}

\subsection{Observables Teleportation:\texorpdfstring{\\}{line}General Framework}
\label{sec:generic_teleportation}

Quantum charge teleportation (QCT) extends Quantum Energy Teleportation (QET)~\cite{Hotta2011,Ikeda2023} from energy to symmetry charges. We consider two closely related $(N\!+\!1)$-site quantum systems prepared in a non-degenerate ground state $\ket{gs}$ of a Hamiltonian $H$, with Alice at site $0$ and Bob at site $N$~\cite{QKDbyQET,IkedaTeleportingCharge}. In QCT, the teleported quantity is a global conserved charge \footnote{The conserved charge can also be discrete, for which there is no continuity equation in the usual Noether sense.}—i.e., the generator of an exact symmetry with an associated continuity equation—so that Alice’s local measurement followed by classical communication (LOCC) enables Bob to apply a conditional local unitary that deterministically activates a charge signal at his site while the total charge remains conserved. This symmetry structure is essential: unlike arbitrary local observables, only globally conserved charges admit nonlocal redistribution via QET-style feedback without intermediate transport. Building on this principle, we analyze how Alice’s measurement basis and Bob’s conditional operation steer the sign and magnitude of Bob’s local charge, thereby realizing charge teleportation suitable for cryptographic primitives based on QET.

\subsubsection{Protocol Steps}
We explicitly define the protocol steps as follows~\cite{QKDbyQET}:

\begin{enumerate}
    \item \textbf{Ground State Preparation:} 
    The global system is initialized in the entangled ground state density matrix $\rho_{gs}=\ketbra{gs}{gs}$.

    \item \textbf{Alice’s Local Measurement:}
    Alice performs a local projective measurement on site $0$ given by:
    \begin{equation}
    P_A(b,\sigma_A)=\frac{1}{2}\left( 1-(-1)^b\sigma_A \right),\quad b\in\{0,1\},
    \end{equation}
    where $\sigma_A=\hat{n}\cdot\vec{\sigma}_A$ is randomly selected from a predefined set (for instance, $\{X_0,Y_0\}$), obtaining measurement outcome $b$.

    \item \textbf{Classical Communication:}
    Alice decides if she communicates her true measurement result ($b$) or its opposite ($b\oplus 1$). This decision is captured by a classical bit $a\in\{0,1\}$, where $a=0$ indicates the true result transmission, and $a=1$ its opposite. Alice communicates the classical bit $c=b\oplus a$ and the chosen measurement basis $\sigma_A$ to Bob. Note that Bob can deduce neither Alice's secret bit $a$ nor her raw measurement outcome $b$ from the public communication of $c$ and $\sigma_A$ alone.

    \item \textbf{Bob’s Conditional Rotation:}
    Based on the classical information received, Bob applies a conditional rotation on his subsystem at site $N$:
    \begin{equation}
    U_B(c,\sigma_B)=e^{-i\theta(-1)^c\sigma_B},
    \end{equation}
    with the operator $\sigma_B$ selected appropriately to match Alice's measurement basis, ensuring maximal teleportation efficiency.

    \item \textbf{Bob’s Observable Measurement:}
    Bob evaluates the expectation shift in his local observable $O_B$ resulting from the protocol:
    \begin{equation}
    \langle\Delta O_B\rangle=\Tr[\rho_B O_B]-\Tr[\rho_{gs} O_B],
    \end{equation}
    where the final density matrix at Bob's site is:
    \begin{align}
    \rho_B=\sum_{b}&{}U_B(b\oplus a,\sigma_B)P_A(b,\sigma_A)\rho_{gs} \nonumber\\
    &P_A(b,\sigma_A)U_B^\dagger(b\oplus a,\sigma_B)
    \end{align}
\end{enumerate}

\subsubsection{Generic Observable Analysis Including Bit \texorpdfstring{$a$}{a}}

Considering a generic observable $O_B$ commuting with Alice’s measurement, i.e., $[P_A,O_B]=0$, we derive explicitly:
\begin{equation}
\label{eq:deltaO-general}
\langle\Delta O_B\rangle=\frac{1}{2}\xi(1-\cos 2\theta)-\frac{1}{2}(-1)^a\eta\sin 2\theta,
\end{equation}
with parameters:
\begin{align}
&{}\xi=\Tr[\rho_{gs}\sigma_B O_B \sigma_B]-\Tr[\rho_{gs}O_B] \nonumber\\
&\eta=i\Tr[\rho_{gs}\sigma_A[O_B,\sigma_B]]
\end{align}

Bob chooses the optimal rotation angle assuming Alice sent her true result ($a=0$), thus maximizing the magnitude of $\langle\Delta O_B\rangle$:
\begin{equation}
\tan(2\theta)=\frac{\eta}{\xi}.
\end{equation}

This leads to the maximal teleported observable expectation shift:
\begin{equation}
\langle\Delta O_B\rangle=\frac{1}{2}\xi-\frac{1}{2}\frac{\xi^2+(-1)^a\eta^2}{\sqrt{\xi^2+\eta^2}}.
\end{equation}

This explicit dependence on $a$ demonstrates the crucial role of Alice's communication decision in secure quantum key distribution (QKD) protocols, and for $a=0$ it is simplified into the known form of:

\begin{equation*}
\langle\Delta O_B\rangle=\frac{1}{2}\xi-\frac{1}{2}\sqrt{\xi^2+\eta^2}.
\end{equation*}

\subsubsection{Complex Sum–Constructed Charge Operator}
\label{sec:complex-sum-constructed-operator}

In what follows we assume, as demanded earlier, that the teleported quantity is a \emph{global conserved charge}
\begin{equation}
[O_B,H]=0,
\end{equation}
with $O_B$ acting at (or supported on) Bob’s site and generated by an exact symmetry of $H$. We further allow $O_B$ to be represented as a complex sum of not-necessarily commuting components, $O_B=\sum_i O_i$, where the $O_i$ are physical observables but, in general, are \emph{not} themselves conserved charges: typically $[O_i,H]\neq 0$, and mutual non-commutativity $[O_i,O_j]\neq 0$ is permitted. Consequently, individual $O_i$ cannot be teleported as charges on their own; rather, the protocol teleports the conserved \emph{sum} $O_B$.

Hence, in order to teleport the charge $O_B$, we analyze each component $O_i$ under this same $\theta$, optimizing $O_B$'s teleportation. The component-wise teleported expectation shift is then
\begin{equation}
\label{eq:deltaOi}
\langle \Delta O_i \rangle
= \frac{1}{2}\,\xi_i
- \frac{1}{2}\,\frac{\xi_i\,\xi + (-1)^a\,\eta_i\,\eta}{\sqrt{\xi^2+\eta^2}},
\end{equation}
with
\begin{align}
\xi_i &= \Tr\!\left[\rho_{gs}\,\sigma_B O_i \sigma_B\right]-\Tr\!\left[\rho_{gs}\,O_i\right], \\
\eta_i &= i\,\Tr\!\left[\rho_{gs}\,[O_i,\sigma_B]\right].
\end{align}
Equation~\eqref{eq:deltaOi} makes explicit that (i) the protocol’s optimization is performed at the level of the conserved $O_B$ (not per $O_i$), and (ii) it reveals the inherent complexity in measuring and analyzing teleported observables composed of non-commuting terms, emphasizing the necessity of separate measurements for each operator component. Such complexity must be carefully considered in practical implementations and QKD protocol designs leveraging charge teleportation~\cite{QKDbyQET,IkedaTeleportingCharge}.

\subsubsection{Discussion and Implications}

The general teleportation protocol for observables discussed herein significantly extends the utility of quantum teleportation beyond energy, demonstrating broad applications for charge, current, and other observables~\cite{IkedaTeleportingCharge}. It lays essential foundations for novel quantum cryptographic techniques, such as QKD schemes, wherein the transmitted classical information (captured by the decision bit $a$) critically influences security and robustness.

Moreover, the presented analysis clarifies the conditions required for successful and optimal teleportation, highlighting the roles of measurement basis selection, classical feedback control, and operator commutation relations. It provides a unified framework facilitating future research and experiments in quantum information, quantum cryptography, and condensed matter physics employing quantum many-body systems~\cite{Hotta2011,Ikeda2023,QKDbyQET}.

\subsection{QKD with Charge teleportation}

We now apply the general observable teleportation protocol to the charge operator. The charge at Bob’s site (site $N$) is defined as:
\begin{equation}
Q_B = \frac{1}{2}(I + Z_N),
\end{equation}
which acts as a projector onto the logical \( \ket{0} \) state. This observable is Hermitian, diagonal in the computational basis, and naturally suited for logical bit interpretation in a QKD scheme.

From the general analysis in Section~\ref{sec:generic_teleportation}, we specialize the result for commuting observables to the case of $Q_B$, yielding:
\begin{equation}
\langle \Delta Q_B \rangle = -\frac{1}{2} \langle Z_N \rangle - \frac{1}{2} \frac{\langle Z_N \rangle^2 + (-1)^a \langle O_A O_B \rangle^2}{\sqrt{\langle Z_N \rangle^2 + \langle O_A O_B \rangle^2}},
\end{equation}
where the classical bit \(a \in \{0,1\}\) encodes whether Alice sent her true measurement result (\(a=0\)) or its complement (\(a=1\)).

The operator pair \((O_A, O_B)\) depends on the chosen measurement and feedback basis:
\begin{equation}
(O_A, O_B) =
\begin{cases}
(X_0, X_N) & \text{for } (\sigma_A, \sigma_B) = (X_0, Y_N), \\
(Y_0, Y_N) & \text{for } (\sigma_A, \sigma_B) = (Y_0, X_N).
\end{cases}
\end{equation}

The teleported expectation value $\langle \Delta Q_B \rangle$ is positive or negative depending on $a$, which enables secure key distribution:
\begin{equation}
\langle \Delta Q_B \rangle
\begin{cases}
< 0 & \Rightarrow \text{logical bit } 1, \\
> 0 & \Rightarrow \text{logical bit } 0.
\end{cases}
\end{equation}

This charge-based observable was specifically chosen for this protocol. As we will demonstrate, its perfectly symmetric response to Alice's classical bit $a$ provides a clean, unambiguous mapping to a logical key bit, forming the basis for a more robust and reliable QKD protocol compared to other observables like energy.

This behavior is at the heart of QKD via observable teleportation: Bob cannot distinguish which logical value was intended unless he knows Alice's classical bit $a$. The security stems from the fact that the expectation value's sign flips based on $a$, yet neither the quantum state nor the classical communicated bit alone is sufficient to decode the logical value.
